\label{ejercicios}
\chapter{Ejercicios y estudios}

\newcounter{exercise}

Este capítulo presenta una serie de ejercicios con tres objetivos: desarrollar la técnica, interiorizar algunos de los patrones rítmicos más importantes de la música árabe, y afianzar el conocimiento de los distintos \maqamaatb.

Se recomienda practicarlos a velocidad moderada, haciendo especial énfasis en lograr una interpretación clara y una buena pulsación, para después ir trabajando a un \emph{tempo} más elevado.

\section{Ejercicios técnicos}

\subsubsection{Ejercicio de \emph{glissando}}

Préstese especial atención a la entonación. La partitura contiene la digitación a emplear.

\audioclip{ejercicio_glissando.mp3}{Ejercicio de \emph{glissando}}

\subsection{Ejercicios de tresillos}

Estos ejercicios están pensados para practicar la pulsación en tresillos según un patrón descendente--ascendente--descendente (\downbow \upbow \downbow). Frases similares a las de esta sección aparecerán en otras partes del capítulo con patrones de pulsación distintos, para trabajar en ese caso esas otras variantes.

\section{Ejercicios rítmicos}

Los siguientes ejercicios se basan en los principales ritmos (\emph{iqaa}) árabes, presentando su estructura principal y algunas variaciones rítimicas.

\subsubsection{Ritmo \emph{malfuf}}

\begin{center}
\begin{lilypond}
    \include "arabic.ly"
    #(set-global-staff-size 18)
    notes = \relative {
      \key do \minor
      \time 3,3,2 2/4
        do8. mib8. mib8
        do8. mib8. mib8
        do8. mib8. mib8
        do8. mib8. mib8
        do8. mib8. mib16 (re)
        do8. mib8. mib8
        do8. mib8. mib16 (re)
        do8. mib8. mib8
        do16-> do16 do16 mib16-> mib16 mib16 mib8->
        do8. mib8. mib8
        do16-> do16 do16 mib16-> mib16 mib16 mib8->
        do8. mib8. mib8
        do16-> do16 do16 mib16-> mib16 mib16 mib16-> (re)
        do8. mib8. mib8
        do16-> do16 do16 mib16-> mib16 mib16 mib16-> (re)
        do8. mib8. mib8
        do16 sol' sol sol lab (sol) fa mib
        do8. mib8. mib8
        do16 sol' sol sol lab (sol) fa mib
        do8. mib8. mib8
        do16 sol' sol sol lab (sol) fa mib
        do16 sol' sol sol lab (sol) fa mib
        do16 sol' sol sol lab (sol) fa mib
        do8. mib8. mib8
        do8. mib8. mib8
        \bar "||"
    }

    \layout {
      \context {
          \TabStaff
        stringTunings = \stringTuning <do, fa, la, re sol do'>
        }
       \context {
          \Staff
          \remove "Time_signature_engraver"
        }
      indent = #0
      \context {
        \Score
        supportNonIntegerFret = ##t
      }
    }

    <<
      \new Staff << \clef "G_8" \notes>>
      \new TabStaff <<\notes>>
    >>
\end{lilypond}
\end{center}

\subsubsection{Ritmo \emph{maqsum}}

\begin{center}
\begin{lilypond}
    \include "arabic.ly"
    #(set-global-staff-size 18)
    notes = \relative {
      \time 4/4
      \key do \minor
      	re8 sol r8 sol fa r8 sol r8
      	re8 sol r8 sol fa fa sol r8
      	re8 sol r8 sol fa r8 sol r8
      	re8 sol r8 sol fa fa sol r8
      	re8 sol sol16 sol sol sol fa8 r8 sol r8
      	re8 sol sol16 sol sol sol fa8 fa sol r8
      	re8 sol sol16 sol sol sol fa8 r8 sol r8
      	re8 sol sol16 sol sol sol la16 (sol) fa8 sol r8
      	re8 sol r8 sol fa r8 sol r8
        re8 la'16 (sib) sol (la) fa (sol) mib (fa) re (mib) do (re) sib (do)
        sol8 sol' r8 sol fa r8 sol r8
        re8 sol r8 sol fa fa sol r8
      	re8 sol r8 sol fa r8 sol r8
        re8 la'16 (sib) sol (la) fa (sol) mib (fa) re (mib) do (re) sib (do)
        sol4 r2.
        \bar "||"
    }

    \layout {
      \context {
          \TabStaff
        stringTunings = \stringTuning <do, fa, la, re sol do'>
        }
       \context {
          \Staff
          \remove "Time_signature_engraver"
        }
      indent = #0
      \context {
        \Score
        supportNonIntegerFret = ##t
      }
    }

    <<
      \new Staff << \clef "G_8" \notes>>
      \new TabStaff <<\notes>>
    >>
\end{lilypond}
\end{center}

\subsubsection{Ritmo \emph{baladi}}

\begin{center}
\begin{lilypond}
    \include "arabic.ly"
    #(set-global-staff-size 18)
    notes = \relative {
      \time 4/4
      \key do \rast
      	re8 re r8 do re r8 do r8
      	re8 re r8 do re r8 do r8
      	re8 re do16 do do8 re do16 do do8 do16 do
      	re8 re do16 do do8 re do16 do do8 do16 do
      	re8 re do16 do do8 re misb16 fa misb re do8
      	re8 re do16 do do8 re misb16 fa misb re do8
      	re8 re do16 do do8 re fa8:64 misb16 re do8
      	re8 re do16 do do8 re fa8:64 misb16 re do8
        \bar "||"
    }

    \layout {
      \context {
          \TabStaff
        stringTunings = \stringTuning <do, fa, la, re sol do'>
        }
       \context {
          \Staff
          \remove "Time_signature_engraver"
        }
      indent = #0
      \context {
        \Score
        supportNonIntegerFret = ##t
      }
    }

    <<
      \new Staff << \clef "G_8" \notes>>
      \new TabStaff <<\notes>>
    >>
\end{lilypond}
\end{center}

\section{Ejercicios por \maqamb}

\subsection{\emph{Maqam} Rast}
\subsection{\emph{Maqam} Bayati}
\subsection{\emph{Maqam} Ayam}
\subsection{\emph{Maqam} Nahawand}

%https://issuu.com/theschoolofarmenianoud/docs/i_girq__lriv